\section{The profile among distinct locations}

Put
\[
H_r(T_1,T_2):=N_d(T_1,T_2)-D_{\leq r}(T_1,T_2),
\]
the number of distinct locations of multiplicity at least $r+1$.

\begin{lemma}[Distinct-location majorant]\label{lem:distmaj}
Fix $r\geq1$.  Suppose $c,\alpha,\beta,\gamma>0$ satisfy $\alpha=2c-\kMT$ and, for every integer $m\geq1$,
\begin{align}
k_c(m)&\leq \alpha m+\beta-\gamma\ind_{m\geq r+1},\label{eq:diston}\\
c^2&\leq 2\alpha m+2\beta-2\gamma\ind_{m\geq r+1}.\label{eq:distpair}
\end{align}
Then
\begin{equation}\label{eq:distgeneric}
\limsup_{T\to\infty}\frac{H_r(T,2T)}{N_d(T,2T)}\leq\frac{\beta}{\gamma}.
\end{equation}
\end{lemma}

\begin{proof}
Summing the pointwise inequalities and using \eqref{eq:countidentity} gives
\[
\sum_m k_c(m)x_m+c^2\sum_m y_m
\leq \alpha N+\beta N_d-\gamma H_r.
\]
For every $\varepsilon>0$, Proposition~\ref{prop:master} and $\alpha=2c-\kMT$ imply, for all sufficiently large $T$,
\[
\gamma H_r\leq\beta N_d+\varepsilon N.
\]
The distinct-zero theorem in \cite{Claude2026} gives $N_d\geq\delta N$ eventually for some fixed $\delta>0$.  Hence
\[
\limsup_{T\to\infty}\frac{H_r}{N_d}
\leq\frac{\beta}{\gamma}+\frac{\varepsilon}{\gamma\delta}.
\]
Let $\varepsilon\downarrow0$.
\end{proof}

\begin{proof}[Proof of Theorem~\ref{thm:distinct}]
For $r=1$, take
\[
c=2,
\quad \alpha=4-\kMT,
\quad \beta=\kMT-1,
\quad \gamma=3-\kMT.
\]
For $m=1$ the on-line inequality is equality.  For $m\geq2$, the right side is increasing in $m$ and equals $k_2(2)=4$ at $m=2$.  The pair inequality has nonnegative slack.  Lemma~\ref{lem:distmaj} gives \eqref{eq:D1}.

For $r=2$, take
\[
c=3,
\quad \alpha=6-\kMT,
\quad \beta=\kMT-1,
\quad \gamma=2(4-\kMT).
\]
Again the on-line inequality is equality at $m=1$ and $m=3$, and has positive slack at $m=2$; the pair inequality is weaker.  This gives \eqref{eq:D2}.

For $r\geq3$, use $C=2+\sqrt2$ and set
\begin{equation}\label{eq:distparams}
\alpha:=2C-\kMT,
\qquad
\beta:=\kMT-1,
\qquad
\gamma:=\alpha(r+1)+\beta-C^2.
\end{equation}
For $m=1$, equation \eqref{eq:balance} gives
\[
k_C(1)=2C-1=\alpha+\beta.
\]
For $m=2,3$ one checks exactly that
\begin{equation}\label{eq:smallslack}
\alpha m+\beta-k_C(m)=(m-1)\bigl(m-(\kMT-1)\bigr)\geq0.
\end{equation}
For $4\leq m\leq r$, $k_C(m)=C^2$.  Since $\alpha>0$, it is enough to check $m=4$; here
\[
4\alpha+\beta-C^2=4C-3\kMT+1>0,
\]
using $\kMT<4/3$.  For $m\geq r+1$, the right side of \eqref{eq:diston} is increasing in $m$ and equals $C^2$ at $m=r+1$ by the definition of $\gamma$.  Thus \eqref{eq:diston} holds.  The pair inequality is equality for a simple pair and has positive slack otherwise.

Lemma~\ref{lem:distmaj} therefore gives
$\limsup_{T\to\infty}H_r(T,2T)/N_d(T,2T)\leq\beta/\gamma$.  Writing $B=(3-2\sqrt2)(\kMT-1)$ and simplifying,
\[
\frac{\beta}{\gamma}=\frac{B}{r-1-rB},
\]
which proves \eqref{eq:Dr}.  Cumulative and Dirichlet variants follow as before.
\end{proof}
